% Full Section 4.4 proofs and the branching prerequisites used by its application.

\begin{annexitem}{annex:p15.bounded-stopping-routes}{thm:p15-bounded-stopping-routes}{Both expectation proofs and the stopped-field refinement}
\begin{proof}
For an almost-surely bounded stopping time \(T\), \(X_T\) agrees
almost surely with \(X_{T\wedge k}\). It is integrable because
\(|X_{T\wedge k}|\leq\sum_{j=0}^k|X_j|\).
The real version specified in the definition is measurable in
\(\mathcal F_T\): intersecting an inverse image with \(\{T\leq n\}\)
gives the finite union over \(j\leq n\) of its
\(\{T=j\}\cap\{X_j\in B\}\) pieces, each in \(\mathcal F_n\).

Stopping preserves the submartingale property, so
\[
                        EX_0=EX_{T\wedge0}
                                  \leq EX_{T\wedge k}=EX_T.
\]
For the source's complementary-transform proof put
\(H_j=\mathbf1_{\{T<j\}}\). This is nonnegative, bounded and
\(\mathcal F_{j-1}\)-measurable. The transform theorem gives
\[
                 E(X_k-X_{T\wedge k})
                   =E\sum_{j=1}^kH_j\Delta X_j\geq0.
\]
This proves \(EX_T\leq EX_k\).
The second source proof instead notes that, for \(j\leq k\),
\[
          E(X_k;T=j)
             =E(E(X_k\mid\mathcal F_j);T=j)
             \geq E(X_j;T=j).
\]
Sum over \(0\leq j\leq k\); the omitted event \(\{T>k\}\)
has probability zero, giving the same conclusion.

For two pointwise ordered times \(S\leq T\), the indicators
\(H_j=\mathbf1_{\{S<j\leq T\}}\) are predictable:
\(\{S<j\},\{T\geq j\}\in\mathcal F_{j-1}\).
Their transform at \(k\) equals
\(X_{T\wedge k}-X_{S\wedge k}\), so \(EX_S\leq EX_T\)
whenever \(P(T\leq k)=1\).

For the gluing assertion,
\[
 \{L\leq n\}=(A\cap\{S\leq n\})\ \cup\ (A^c\cap\{T\leq n\}).
\]
The first piece is in \(\mathcal F_n\) by \(A\in\mathcal F_S\).
Pointwise ordering gives \(\{T\leq n\}\subseteq\{S\leq n\}\);
the second piece is therefore
\(\{T\leq n\}\setminus(A\cap\{S\leq n\})\), also measurable there.
Thus \(L\) is a stopping time and \(L\leq T\).
Under the bound, the preceding expectation inequality yields
\[
                 E(X_S;A)+E(X_T;A^c)=EX_L\leq EX_T.
\]
Hence \(E(X_S;A)\leq E(X_T;A)\) for every \(A\in\mathcal F_S\).
Both variables are integrable and \(X_S\) is
\(\mathcal F_S\)-measurable. The defining conditional identity and
the test sets where \(X_S-E(X_T\mid\mathcal F_S)>1/r\)
prove \(X_S\leq E(X_T\mid\mathcal F_S)\) almost surely.
Changing sign gives the supermartingale statement; applying both
inequalities to a martingale gives equality.
\end{proof}
\end{annexitem}

\begin{annexitem}{annex:p15.doob-weak-maximal}{thm:p15-doob-weak-maximal}{Stopping at the first threshold crossing}
\begin{proof}
Let \(T=\inf\{j\geq0:X_j\geq\lambda\}\wedge n\).
The threshold event up to any time is a finite union of adapted
events, so this is a stopping time bounded by \(n\).
On \(A=\{\overline X_n\geq\lambda\}\), one has \(X_T\geq\lambda\).
On \(A^c\), no threshold was reached and \(X_T=X_n\).
The bounded expectation inequality \(EX_T\leq EX_n\), after
subtracting the equal contributions on \(A^c\), gives
\[
                      \lambda P(A)\leq E(X_T;A)\leq E(X_n;A).
\]
Finally \(X_n\mathbf1_A\leq X_n^+\) pointwise.
For a strict crossing replace the threshold condition by
\(X_j>\lambda\); the same proof gives the stated weak bound
for its strict event.
The process \(|Y_j|\) is an integrable submartingale by conditional
Jensen for the convex absolute-value function. Applying the first
result to it proves the absolute-maximum assertion.
\end{proof}
\end{annexitem}

\begin{annexitem}{annex:p15.kolmogorov-martingale-route}{thm:p15-kolmogorov-martingale-route}{Squaring the independent-sum martingale}
\begin{proof}
In the filtration \(\sigma(\xi_1,\ldots,\xi_m)\), with trivial
time-zero field, independence and centering give
\(E(S_{m+1}\mid\mathcal F_m)=S_m\).
Finite variance addition gives \(ES_n^2=v_n\).
The square process is integrable and is a submartingale by
conditional Jensen. Its positive maximum is
\(\max_{m\leq n}S_m^2\). Doob's inequality at threshold \(x^2>0\)
therefore gives
\[
             x^2P(\max_{m\leq n}|S_m|\geq x)\leq ES_n^2=v_n.
\]
Division by \(x^2\) proves the result, including \(v_n=0\).
The deterministic time-zero term is zero and does not change the event.
\end{proof}
\end{annexitem}

\begin{annexitem}{annex:p15.doob-lp-maximal}{thm:p15-doob-lp-maximal}{The complete truncated layer-cake and Holder argument}
\begin{proof}
If \(E(X_n^+)^p=\infty\), the claimed extended inequality is immediate.
Otherwise fix \(M>0\) and put \(Z=\overline X_n\wedge M\).
The power tail identity, Doob's inequality, and Tonelli for the
nonnegative integrand give
\[
 \begin{aligned}
 EZ^p
   &=\int_0^M p\lambda^{p-1}P(\overline X_n>\lambda)\,d\lambda\\
   &\leq\int_0^M p\lambda^{p-2}
                          E(X_n^+;\overline X_n>\lambda)\,d\lambda\\
   &=E\left[X_n^+\int_0^Zp\lambda^{p-2}\,d\lambda\right]
     =\frac p{p-1}E(X_n^+Z^{p-1}).
 \end{aligned}
\]
The integral at zero is finite since \(p>1\).
With \(q=p/(p-1)\), Holder's inequality bounds the last expression by
\[
                    \frac p{p-1}
                         \bigl(E(X_n^+)^p\bigr)^{1/p}
                         (EZ^p)^{(p-1)/p}.
\]
The number \(EZ^p\) is finite because \(Z\leq M\).
If it is zero, the desired truncated bound already holds.
Otherwise divide by its positive \((p-1)/p\) power and raise
to power \(p\). The resulting inequality is
\[
                         EZ^p\leq
                              \left(\frac p{p-1}\right)^pE(X_n^+)^p.
\]
Monotone convergence as \(M\to\infty\) proves the first result.
Apply it to the submartingale \(|Y|\) for the second.
\end{proof}
\end{annexitem}

\begin{annexitem}{annex:p15.fair-walk-maximal-tail}{thm:p15-fair-walk-maximal-tail}{Finite-interval exit and the harmonic maximal tail}
\begin{proof}
The walk is an integrable martingale in its independent-sign
filtration, and stopping at \(T_0\) gives a martingale with mean one
at every finite time. Unit increments cannot cross below zero
without first reaching it, so the stopped process is nonnegative.

Fix an integer \(m\geq2\), let \(T_m\) be the first hit of \(m\),
and put \(T=T_0\wedge T_m\).
Before this exit, the walk is at one of \(1,\ldots,m-1\).
Given the information at time \(jm\), if exit has not occurred,
the event that all of the next \(m\) increments are \(-1\)
has conditional probability \(2^{-m}\) and forces exit at zero.
Induction over these blocks gives
\[
                    P(T>jm)\leq(1-2^{-m})^j .
\]
Thus \(T<\infty\) almost surely.
For every finite \(n\), bounded optional sampling gives
\(ES_{T\wedge n}=1\), and \(0\leq S_{T\wedge n}\leq m\).
Bounded convergence yields \(ES_T=1\).
The terminal value is \(m\) on \(\{T_m<T_0\}\) and zero otherwise;
simultaneous hits are impossible. Hence \(P(T_m<T_0)=1/m\).
For \(m=1\), the initial value is already one, so the formula holds too.

On \(\{T_0=\infty\}\), the finite exit \(T_0\wedge T_m\)
must be \(T_m<T_0\). Therefore \(P(T_0=\infty)\leq1/m\)
for every \(m\geq2\), and \(T_0<\infty\) almost surely.
Reaching height at least \(m\) before absorption is equivalent to
first hitting \(m\), because the increments are unit steps.
This proves the claimed maximal tail.

The all-time maximum is a finite integer almost surely,
because absorption occurs at a finite time. For an integer-valued
nonnegative variable, Tonelli gives its expectation as the sum
of its positive integer tail probabilities. Consequently
\[
                  E\sup_nX_n=\sum_{m=1}^\infty\frac1m=\infty.
\]
The harmonic series diverges, for example by grouping consecutive
dyadic blocks, each of sum at least \(1/2\).
Finite-horizon maxima increase to this supremum, so their
expectations tend to infinity by monotone convergence.
Since every terminal expectation remains one, no finite constant
can give the asserted universal \(L^1\) maximal bound.
\end{proof}
\end{annexitem}

\begin{annexitem}{annex:p15.martingale-lp-convergence}{thm:p15-martingale-lp-convergence}{An integrable dominating maximum}
\begin{proof}
Holder's inequality on a probability space gives
\(E|X_n|\leq(E|X_n|^p)^{1/p}\leq C^{1/p}\).
Thus the positive-part first moments are bounded, and the
martingale convergence theorem gives a finite almost-sure limit.
The finite maximal inequality and monotone convergence in \(n\)
give
\[
                        E(X_\infty^*)^p
                              \leq\left(\frac p{p-1}\right)^pC<\infty.
\]
On the common convergence event, \(|X_\infty|\leq X_\infty^*\).
Choose a finite measurable limit version, equal to zero off that
event. Then
\[
                    |X_n-X_\infty|^p\leq(2X_\infty^*)^p
                                    \quad\text{almost surely}.
\]
Dominated convergence proves \(L^p\) convergence.
Fatou gives \(E|X_\infty|^p\leq C\).
Holder also gives \(L^1\) convergence from \(L^p\) convergence;
the closed-martingale theorem therefore identifies
\(X_n=E(X_\infty\mid\mathcal F_n)\).

For a nonnegative submartingale with bounded \(p\)th moments,
the same first-moment estimate gives almost-sure convergence.
Its absolute maximum equals its positive maximum, so the first
maximal inequality supplies the identical integrable dominating
variable. The same dominated-convergence proof gives \(L^p\)
convergence. The subsequent counterexamples settle the claimed
limitations.
\end{proof}
\end{annexitem}

\begin{annexitem}{annex:p15.martingale-orthogonality}{thm:p15-martingale-orthogonality}{Earlier information and every cross term}
\begin{proof}
Cauchy--Schwarz gives \(E|(X_n-X_m)Y|<\infty\).
The conditional pullout rule for this integrable product, followed
by the martingale tower identity, gives
\[
 E((X_n-X_m)Y)
      =E\bigl[Y\,E(X_n-X_m\mid\mathcal F_m)\bigr]=0.
\]
For \(i<j\), use \(X_j-X_{j-1}\) and the
\(\mathcal F_{j-1}\)-measurable \(L^2\) variable \(\Delta X_i\);
their product has expectation zero.
The same holds with \(X_0\) in place of \(\Delta X_i\).
Expanding the square of
\(X_n-X_0=\sum_{j=1}^n\Delta X_j\) and dropping precisely these
zero cross expectations gives the first identity.
Expanding \(X_n=X_0+\sum_j\Delta X_j\) gives the second.
Every sum here is finite and every product is integrable by
Cauchy--Schwarz.
\end{proof}
\end{annexitem}

\begin{annexitem}{annex:p15.martingale-conditional-variance}{thm:p15-martingale-conditional-variance}{The square expansion and its projection instance}
\begin{proof}
The square and cross terms are integrable, so conditional linearity
and the pullout rule are applicable:
\[
 \begin{aligned}
 E((X_n-X_m)^2\mid\mathcal F_m)
   &=E(X_n^2\mid\mathcal F_m)
       -2X_mE(X_n\mid\mathcal F_m)+X_m^2\\
   &=E(X_n^2\mid\mathcal F_m)-X_m^2.
 \end{aligned}
\]
For the second assertion, put \(Z_0=E(Y\mid\mathcal H)\),
\(Z_1=E(Y\mid\mathcal G)\).
Conditional \(L^2\) contraction makes both square integrable, and
the tower identity gives \(E(Z_1\mid\mathcal H)=Z_0\).
They form a two-time martingale in the nested fields.
Applying the first formula and taking expectations gives
\(E(Z_1-Z_0)^2=EZ_1^2-EZ_0^2\), exactly the asserted identity.
\end{proof}
\end{annexitem}

\begin{annexitem}{annex:p15.martingale-covariance-isometry}{thm:p15-martingale-covariance-isometry}{Telescoping covariance and transformed increments}
\begin{proof}
At each time,
\[
 X_jY_j-X_{j-1}Y_{j-1}
      =X_{j-1}\Delta Y_j+Y_{j-1}\Delta X_j+\Delta X_j\Delta Y_j.
\]
All products are integrable. The first two expectations vanish
by orthogonality to the common earlier information.
Take expectations and sum the finite telescoping identity to
obtain the covariance formula.

The predictable transforms
\(\widetilde X_n=\sum_{j=1}^nH_j\Delta X_j\) and
\(\widetilde Y_n=\sum_{j=1}^nK_j\Delta Y_j\)
are martingales by the transform theorem.
Their finite-time values are in \(L^2\): each bounded coefficient
times an \(L^2\) increment is in \(L^2\), and there are finitely
many terms. They start at zero and their increments are
\(H_j\Delta X_j\) and \(K_j\Delta Y_j\).
Apply the first covariance formula to these two transforms.
Taking \(Y=X\), \(K=H\) gives the square isometry.
\end{proof}
\end{annexitem}

\begin{annexitem}{annex:p15.square-summable-martingale}{thm:p15-square-summable-martingale}{The limit and the exact tail error}
\begin{proof}
The moment hypothesis makes every finite
\(X_n=X_0+\sum_{j\leq n}\Delta X_j\) square integrable.
Orthogonality gives
\[
              EX_n^2=EX_0^2+\sum_{j=1}^nE(\Delta X_j)^2
                       \leq EX_0^2+\sum_{j=1}^\infty E(\Delta X_j)^2.
\]
The \(p=2\) convergence theorem supplies a limit, almost surely
and in \(L^2\).
For fixed \(n\) and \(m>n\), orthogonality also gives
\[
                         E(X_m-X_n)^2
                                   =\sum_{j=n+1}^mE(\Delta X_j)^2.
\]
The \(L^2\) triangle inequality shows that
\(\|X_m-X_n\|_2\to\|X_\infty-X_n\|_2\).
Squaring this limit and passing the nonnegative finite sums
to their increasing limit proves the error formula.
The last assertion applies this argument to the martingale
of partial sums of the specified martingale differences.
\end{proof}
\end{annexitem}

\begin{annexitem}{annex:p15.weighted-martingale-law}{thm:p15-weighted-martingale-law}{A martingale series followed by Kronecker's lemma}
\begin{proof}
The deterministic coefficients \(1/b_j\) are finite and predictable.
Thus
\[
                         Y_n=\sum_{j=1}^n\frac{\Delta X_j}{b_j}
\]
is a square-integrable martingale starting at zero.
The sum of the second moments of its increments is finite by
hypothesis. The square-summable theorem makes this series converge
almost surely to a finite limit.
On its convergence event the numerical Kronecker lemma, applied
to \(x_j=\Delta X_j\) with the specified positive nondecreasing
weights \(b_j\), gives
\[
                  \frac{X_n-X_0}{b_n}
                         =\frac1{b_n}\sum_{j=1}^n\Delta X_j\longrightarrow0.
\]
Since \(X_0\) is finite almost surely and \(b_n\to\infty\),
also \(X_0/b_n\to0\).
If \(E(\Delta X_j)^2\leq K\), the required series is at most
\(K\sum_jb_j^{-2}<\infty\), proving the stated sufficient condition.
\end{proof}
\end{annexitem}

\begin{annexitem}{annex:p15.bounded-increment-small-ball}{thm:p15-bounded-increment-small-ball}{The strict exit and its bounded overshoot}
\begin{proof}
The independent-sum conditional square identity shows that
\(Q_j=S_j^2-v_j\) is a martingale with \(Q_0=0\).
Let \(T=\inf\{j\geq1:|S_j|>x\}\wedge n\).
This is bounded by \(n\). If the exit occurs, its preceding
value has absolute value at most \(x\), so \(|S_T|\leq x+K\).
If no exit occurs before the horizon, \(|S_T|\leq x\);
the same upper bound follows.
All these assertions hold outside the common null set for
the finitely many increment bounds.
Bounded optional sampling for \(Q\) gives
\[
                              ES_T^2=Ev_T\leq(x+K)^2.
\]
On \(A=\{\max_{1\leq j\leq n}|S_j|\leq x\}\) one has \(T=n\).
Since \(v_j\geq0\), \(Ev_T\geq v_nP(A)\).
Division by \(v_n>0\) proves the assertion.
\end{proof}
\end{annexitem}

\begin{annexitem}{annex:p15.cantelli-maximal}{thm:p15-cantelli-maximal}{The optimized shift and an attaining law}
\begin{proof}
For any \(c\geq0\), \((X_j+c)^2\) is an integrable submartingale.
On the event that some \(X_j\geq\lambda\), its maximum is at least
\((\lambda+c)^2\). Doob's inequality therefore gives
\[
 P(\max_{j\leq n}X_j\geq\lambda)
          \leq\frac{E(X_n+c)^2}{(\lambda+c)^2}
          =\frac{v+c^2}{(\lambda+c)^2},
\]
since \(EX_n=EX_0=0\).
At \(c=v/\lambda\), direct cancellation gives
\(v/(v+\lambda^2)\); this choice also gives zero when \(v=0\).
To see that the value cannot be lowered uniformly, for \(v>0\)
let a single centered increment take value \(\lambda\) with
probability \(v/(v+\lambda^2)\) and value \(-v/\lambda\)
with the remaining probability \(\lambda^2/(v+\lambda^2)\).
Its mean is zero, and its second moment is
\[
 \frac{v\lambda^2}{v+\lambda^2}
        +\frac{\lambda^2}{v+\lambda^2}\frac{v^2}{\lambda^2}=v.
\]
Starting at zero and holding this value after time one gives
a martingale attaining the displayed crossing probability.
\end{proof}
\end{annexitem}

\begin{annexitem}{annex:p15.doob-llogl}{thm:p15-doob-llogl}{The logarithmic endpoint estimate with all zero cases}
\begin{proof}
If \(E(X_n^+\log^+X_n^+)=\infty\), the bound is immediate.
Otherwise put \(a=X_n^+\) and \(Z=\overline X_n\wedge M\), \(M\geq1\).
The tail integral up to one is at most one.
Above one, use the strict version of Doob's inequality and Tonelli:
\[
 \begin{aligned}
 EZ&=\int_0^MP(\overline X_n>t)\,dt\\
   &\leq1+\int_1^M t^{-1}E(a;\overline X_n>t)\,dt
    =1+E(a\log^+Z).
 \end{aligned}
\]
Define \(\log^+0=0\) in this expression.
For any nonnegative numbers \(a,b\),
\[
                           a\log^+b\leq a\log^+a+b/e,
\]
with \(a\log^+a=0\) at zero. If \(b\leq1\) or \(a=0\),
the assertion is immediate. Otherwise put \(u=a/b>0\);
calculus gives \(u\log(1/u)\leq1/e\), since its derivative
vanishes only at \(u=e^{-1}\) and its maximum there is \(1/e\).
Thus
\(a\log b=a\log a+b\,u\log(1/u)\leq a\log^+a+b/e\).
Apply this pointwise to the random \(a,Z\) above:
\[
                         EZ\leq1+E(a\log^+a)+e^{-1}EZ.
\]
The truncated expectation is finite, so subtraction is legitimate.
Rearrange and let \(M\to\infty\) by monotone convergence.
Finally apply the result to the submartingale \(X/a_0\) for
any deterministic \(a_0>0\), multiply by \(a_0\), and rename
that scale \(a\). This gives the second bound.
\end{proof}
\end{annexitem}

\begin{annexitem}{annex:p15.branching-family-construction}{thm:p15-branching-family-construction}{The countable genealogy and its source-law bridge}
\begin{proof}
The set of finite positive-integer sequences is countable, so the
proved countable independent-product construction supplies all
\(\xi_v\). Induction shows that \(V_n\) is a finite set of
length-\(n\) sequences determined by the offspring values of smaller
length. Thus \(Z_n\) is finite and \(\mathcal H_n\)-measurable.
The still unused length-\(n\) offspring variables are independent
of \(\mathcal H_n\).

There are countably many possible finite sets \(V_n\).
On each event \(\{V_n=A\}\), the conditional law of
\(\sum_{v\in A}\xi_v\) is the law of a sum of
\(\operatorname{card}(A)\) independent offspring variables.
It follows that the conditional distribution of \(Z_{n+1}\)
depends on the past only through \(Z_n\), with transition
\[
              K(k,\{\ell\})=P(\xi_1+\cdots+\xi_k=\ell),
\]
and \(K(0,\{0\})=1\). These are exactly the transitions of the
source's independent-row recursion. The common initial value
and induction on finite histories give equal finite-dimensional
generation-count laws; uniqueness on the countable coordinate
sigma-algebra gives equal process laws.

The nonnegative conditional-mean calculation gives
\(E(Z_{n+1}\mid\mathcal H_n)=\mu Z_n\).
Starting from \(EZ_0=1\), taking expectations and inducting gives
\(EZ_n=\mu^n<\infty\) for \(n\geq1\). Thus these identities
are ordinary integrable conditional expectations at every time.
For \(\mu>0\), division by \(\mu^{n+1}\) proves
\(E(W_{n+1}\mid\mathcal H_n)=W_n\) and \(EW_n=1\).
This is a nonnegative martingale, so it has a finite almost-sure
limit by the already proved convergence theorem.
In the source's row filtration, the next entire row is independent
of all rows already revealed. The same conditional sum calculation
therefore gives the same martingale identity for its normalized process,
including the time-zero step.

For each potential first child \(i\geq1\), the variables
\(\{\xi_{iv}:v\in\mathcal T\}\) form an independent family with
the original law. These countable families are mutually
independent and independent of \(\xi_\varnothing\), by generated
independence of disjoint coordinate sets.
Their realized subtree processes \(Z_n^{(i)}\) are therefore
independent copies of \(Z_n\). Directly from the genealogy,
\[
                    Z_{n+1}=\sum_{i=1}^{\xi_\varnothing}Z_n^{(i)}.
\]
The same construction works below any realized generation, since
its descendant-coordinate families are disjoint and have not
been revealed in \(\mathcal H_n\). This proves the family
independence needed in the subsequent arguments.
\end{proof}
\end{annexitem}

\begin{annexitem}{annex:p15.branching-extinction-root}{thm:p15-branching-extinction-root}{Iteration, critical extinction and the full root argument}
\begin{proof}
Absorption at zero makes \(\{Z_n=0\}\) an increasing sequence of
events. The independent-family decomposition, conditioned on
\(\xi_\varnothing=k\), gives
\[
            P(Z_{n+1}=0\mid\xi_\varnothing=k)=\theta_n^k,
                         \qquad\theta_0=0.
\]
For \(k=0\) the empty family is extinct and the empty product is one.
Summing over \(k\) gives \(\theta_{n+1}=f(\theta_n)\).
The series for \(f\) converges uniformly on \([0,1]\), because
its nonnegative tail is bounded by \(\sum_{k>M}p_k\).
Thus \(f\) is continuous and nondecreasing, with \(f(1)=1\).
Continuity of probability gives \(\theta_n\uparrow\rho=P(D)\);
continuity of \(f\) then yields \(\rho=f(\rho)\).
For any fixed point \(r\in[0,1]\), induction gives
\(\theta_0=0\leq r\) and
\(\theta_{n+1}=f(\theta_n)\leq f(r)=r\).
Taking limits proves that \(\rho\) is the smallest fixed point.

If \(\mu=0\), nonnegativity of the offspring variable forces it
to be zero almost surely, so \(Z_1=0\).
If \(0<\mu<1\), then
\[
                            P(Z_n>0)\leq EZ_n=\mu^n.
\]
The sum of these bounds is finite. Borel--Cantelli gives only
finitely many nonzero generations, hence extinction.

Suppose next that \(\mu=1\).
Then \(Z_n\) is a nonnegative martingale and converges almost surely
to a finite limit. An integer-valued convergent sequence is
eventually constant: once it is within an interval of length
less than one around its limit, it can take at most one integer value.
If \(p_1<1\), then \(p_0>0\). Otherwise offspring would be at
least one almost surely, and mean one would force it to equal one.
For each \(k\geq1\), given \(Z_n=k\), all \(k\) offspring counts
are zero with probability \(p_0^k>0\).
For fixed \(N,k\), successive conditioning bounds
\[
 P(Z_N=k,Z_{N+1}=k,\ldots,Z_{N+r}=k)
                                 \leq(1-p_0^k)^r .
\]
Letting \(r\to\infty\) gives probability zero for remaining at
that positive \(k\) forever. A countable union over \(N,k\)
shows that the eventual constant can only be zero.
If \(p_1=1\), direct induction instead gives \(Z_n=1\) at every time.

For the supercritical root argument, the source differentiations
can be justified directly. On any interval \([0,r]\), \(r<1\),
the derivative series
\[
             \sum_{k\geq1}kp_ks^{k-1},\qquad
             \sum_{k\geq2}k(k-1)p_ks^{k-2}
\]
converge uniformly. Indeed they are dominated by convergent
series of polynomial factors times powers of \(r\).
For example, their successive positive terms without \(p_k\)
have ratio tending to \(r<1\), so beyond a finite index they
are bounded by a geometric tail with ratio strictly below one.
The finite polynomial identities
\(f_M(s)=p_0+\int_0^sf_M'(u)\,du\) and
\(f_M'(s)=p_1+\int_0^sf_M''(u)\,du\)
pass to their uniform limits. The fundamental theorem of
calculus therefore yields
\[
       f'(s)=\sum_{k\geq1}kp_ks^{k-1},\qquad
       f''(s)=\sum_{k\geq2}k(k-1)p_ks^{k-2}\quad(0<s<1).
\]
Monotone convergence as \(s\uparrow1\) gives \(f'(s)\uparrow\mu\).
Since \(\mu<\infty\), integrate this derivative to the endpoint:
\[
        \frac{f(1)-f(1-h)}h
                   =\frac1h\int_{1-h}^1 f'(s)\,ds\longrightarrow\mu.
\]
The endpoint integral follows first below one and then by bounded
convergence, using \(0\leq f'(s)\leq\mu\).
If \(\mu>1\), it follows that \(f(1-h)<1-h\) for small \(h>0\).
Since \(f(0)\geq0\), the intermediate-value theorem gives a fixed
point below one, so the smallest fixed point \(\rho\) is below one.
Also some \(p_k>0\) with \(k\geq2\); the term \(p_ks^k\)
makes \(f\) strictly convex on \([0,1]\).
Between the roots \(\rho\) and 1, strict convexity gives
\(f(s)<s\). There is no root below the smallest root and none
between these two. Hence \(\rho\) is the unique root in \([0,1)\).
\end{proof}
\end{annexitem}

\begin{annexitem}{annex:p15.branching-second-moment}{thm:p15-branching-second-moment}{Every conditional variance and geometric sum}
\begin{proof}
Given the revealed generation \(V_{n-1}\), the next offspring
variables are independent with mean \(\mu\) and variance \(\sigma^2\).
On \(\{Z_{n-1}=k\}\), their sum consequently has mean \(\mu k\)
and variance \(\sigma^2k\). Thus
\[
 \begin{split}
 E(Z_n^2\mid\mathcal H_{n-1})
                   &=\mu^2Z_{n-1}^2+\sigma^2Z_{n-1},\\
 E((Z_n-\mu Z_{n-1})^2\mid\mathcal H_{n-1})
                   &=\sigma^2Z_{n-1}.
 \end{split}
\]
Starting from \(Z_0=1\), the first identity proves by induction
that \(EZ_n^2<\infty\) at every finite time; hence all subsequent
conditional variance operations are legitimate.
Dividing the second identity by \(\mu^{2n}\) gives the claimed
conditional increment formula for \(W_n=Z_n/\mu^n\).

Conditional variance for the martingale \(W\) and then expectation give
\[
          EW_n^2=EW_{n-1}^2+\sigma^2\mu^{-2n}EZ_{n-1}
                     =EW_{n-1}^2+\sigma^2\mu^{-n-1}.
\]
Here \(EZ_{n-1}=\mu^{n-1}\), including \(n=1\).
Since \(EW_0^2=1\), finite summation yields
\[
       EW_n^2=1+\sigma^2\sum_{j=1}^n\mu^{-j-1}
                  =1+\frac{\sigma^2}{\mu(\mu-1)}(1-\mu^{-n}).
\]
The last equality is the finite geometric-sum formula and uses
\(\mu>1\).
The moments are uniformly bounded, so the \(p=2\) convergence
theorem gives the almost-sure and \(L^2\) limit.
The inequality \(\|W_n-W\|_1\leq\|W_n-W\|_2\)
shows convergence of expectations, giving \(EW=1\).
The \(L^2\) triangle inequality gives
\(\|W_n\|_2\to\|W\|_2\), yielding the displayed second moment.
Since \(W\geq0\) and has mean one, it cannot vanish almost surely.
\end{proof}
\end{annexitem}

\begin{annexitem}{annex:p15.branching-positive-limit}{thm:p15-branching-positive-limit}{The limit fixed point and both extinction arguments}
\begin{proof}
For every potential first child \(i\), let
\(W_n^{(i)}=Z_n^{(i)}/\mu^n\).
The independent-family construction and nonnegative martingale
convergence give limits \(W^{(i)}\) with the same law as \(W\).
The countably many family limits and the root limit hold together
on one event of probability one. On that event the number
\(\xi_\varnothing\) of realized first children is finite, so the
finite family identity passes to the limit:
\[
                            W=\frac1\mu
                                      \sum_{i=1}^{\xi_\varnothing}W^{(i)}.
\]
The limits \(W^{(i)}\) are independent of each other and of
\(\xi_\varnothing\), because they are measurable functions of
the corresponding independent descendant-coordinate families.

Put \(r=P(W=0)\).
All terms in the sum are nonnegative, so, conditional on
\(\xi_\varnothing=k\), the sum is zero exactly when all \(k\)
limits are zero. Independence gives probability \(r^k\),
with empty product one. Thus \(r=f(r)\).
The hypothesis that \(W\) is not identically zero gives \(r<1\).
The unique fixed point below one is \(\rho\), so \(r=\rho\).
Extinction forces \(W_n=0\) eventually and hence \(W=0\).
These two events have the same probability \(\rho\), so they
agree up to a null set. This proves the asserted survival equality.

With \(z\) independent initial trees, extinction occurs exactly when
all \(z\) trees become extinct, with probability \(\rho^z\);
for \(z=0\) the population is already zero and the probability is one.
The source's separate bounded-martingale derivation is as follows.
Conditional on a generation of size \(k\), its next-size
probability generating function at \(\rho\) is \(f(\rho)^k=\rho^k\).
Thus \(\rho^{Z_n}\), with \(0^0=1\) in this generating-function
convention, is a bounded martingale.
If \(0<\rho<1\), then \(p_0>0\); otherwise zero itself would
be the smallest fixed point.
The already proved extinction-or-escape theorem
\cref{thm:p15-extinction-or-escape} gives eventual zero or
\(Z_n\to\infty\), almost surely.
Therefore \(\rho^{Z_n}\to\mathbf1_D\), and bounded convergence
of its constant expectation \(\rho^z\) gives \(P(D)=\rho^z\).
If \(\rho=0\), then \(p_0=0\), so a positive initial population
can never become zero; the formula again gives zero for \(z>0\)
and one for \(z=0\).

More explicitly, given any realized generation, its finitely many
descendant families use disjoint unrevealed coordinate sets.
Conditioning on each possible realized finite vertex set and using
their independence gives
\(P(D\mid\mathcal H_n)=\rho^{Z_n}\).
This also identifies the bounded martingale as the current extinction
probability, including the already extinct generation.
\end{proof}
\end{annexitem}

\begin{annexitem}{annex:p15.branching-calculations}{thm:p15-branching-calculations}{The exact binomial root and the degenerate regimes}
\begin{proof}
For a sum of three independent Bernoulli \((1/2)\) offspring
indicators, the mean is \(3/2\), variance \(3/4\), and
\[
                           f(s)=\frac{1+3s+3s^2+s^3}{8}.
\]
The equation \(f(s)=s\) is
\[
               s^3+3s^2-5s+1=(s-1)(s^2+4s-1)=0.
\]
Its unique root in \([0,1)\) is \(\sqrt5-2\).
The preceding survival theorem gives
\(P(W>0)=1-\rho=3-\sqrt5\).
The second-moment formula gives
\[
                         EW^2=1+\frac{3/4}{(3/2)(1/2)}=2.
\]
The generating-function iteration starts with
\(\theta_0=0\), \(\theta_1=1/8\), and
\(\theta_2=((1+1/8)/2)^3=729/4096\);
the extinction-root theorem proves monotone convergence of all
its subsequent iterates to \(\sqrt5-2\).

For deterministic offspring \(b\geq2\), direct induction gives
\(Z_n=b^n\), so \(W_n=1\) and extinction probability is zero.
For binomial \((2,1/2)\) offspring, \(\mu=1\) and \(p_1=1/2<1\).
Critical extinction gives \(Z_n\to0\) almost surely, while the
mean identity gives \(EZ_n=1\) for all \(n\).
Since \(Z_n\geq0\), its \(L^1\) distance from zero remains one.
This proves the asserted failure of convergence in mean.
\end{proof}
\end{annexitem}

\begin{annexitem}{annex:p15.maximal-convergence-boundaries}{thm:p15-maximal-convergence-boundaries}{Every conditional and norm calculation in the boundary models}
\begin{proof}
For the first construction write
\(Y_n=-X_n=(n+1)^{1/p}\mathbf1_{A_n}\).
The event \(A_n\) is an atom of \(\mathcal F_n\);
on its complement all subsequent \(Y\)'s vanish.
On \(A_n\),
\[
 E(Y_{n+1}\mid\mathcal F_n)
        =(n+2)^{1/p}\frac{P(A_{n+1})}{P(A_n)}
        =(n+1)(n+2)^{1/p-1}
        \leq(n+1)^{1/p},
\]
because \(1/p-1<0\).
Thus \(Y\) is a nonnegative supermartingale and \(X=-Y\)
is a submartingale. Each moment is
\(E|X_n|^p=(n+1)P(A_n)=1\).
For every outcome \(\omega>0\), eventually \(\omega>1/(n+1)\);
therefore \(X_n(\omega)=0\) eventually.
The \(p\)th-moment distance to this limit is nonetheless one.

The finite maximum \((X_n^*)^p\) is an integer in
\(\{1,\ldots,n+1\}\).
For \(1\leq k\leq n+1\), its event of being at least \(k\)
is exactly \(A_{k-1}\), since the \(A_j\) are nested.
The finite integer tail identity gives
\[
                         E(X_n^*)^p=\sum_{k=1}^{n+1}P(A_{k-1})
                                             =\sum_{k=1}^{n+1}\frac1k.
\]

For the second construction the dyadic conditional-average theorem
gives almost-sure and \(L^1\) convergence to the indicator.
The bound \(0\leq X_n\leq1\) and the finite-\(p\) convergence
theorem give every \(1<p<\infty\) as well.
At level \(n\), all cells except the one containing \(1/3\)
have a constant indicator and hence zero approximation error.
Writing \(k_n=\lfloor2^n/3\rfloor\), the fraction of the cut
cell to the left of \(1/3\) is
\(\alpha_n=2^n/3-k_n\).
Since powers of two have residue 1 or 2 modulo 3,
\(\alpha_n\in\{1/3,2/3\}\).
The conditional average on that cell is \(\alpha_n\).
Both its indicator-one and indicator-zero portions have positive
length, and their errors are \(1-\alpha_n\) and \(\alpha_n\).
Thus the essential supremum is exactly \(2/3\) at every level.

Finally, \(E\xi_1=0\), and independence of the fair signs gives
\[
                   E(\xi_2\mid\varepsilon_1)
                           =(1+\varepsilon_1)E\varepsilon_2=0.
\]
Their partial sums are a square-integrable martingale in the
input filtration, held constant after the second step.
The same conditional calculation gives
\(E(\xi_1\xi_2)=0\).
But \(\xi_2=0\) exactly when \(\varepsilon_1=-1\), so its
conditional probability there is one and its unconditional
probability is \(1/2\). This proves dependence despite orthogonality.
\end{proof}
\end{annexitem}
